\documentclass[a4paper]{article}
\usepackage[fontset=fandol]{ctex}
\usepackage[margin=2cm]{geometry}
\usepackage{pdfpages}
\usepackage{hyperref}
\usepackage{longtable}
\usepackage{enumitem}
\title{CS 294-277 / S294-277 Robots That Learn（Fall 2024）全课程讲义}
\author{Codex harness-managed textbook build}
\date{\today}
\begin{document}
\maketitle
\tableofcontents
\newpage
\section*{课程说明}
本讲义采用 harness-managed 流水线逐讲生成，并在通过 evaluator 与 validator 后合并。
\addcontentsline{toc}{section}{课程说明}
\section*{官方视频源说明}
\addcontentsline{toc}{section}{官方视频源说明}
本次重建以官方 YouTube 频道与其公开 playlist 为主视频索引，而不是仅依据课程页课表标题做推断。
官方频道：\href{https://www.youtube.com/@robots-that-learn}{\nolinkurl{https://www.youtube.com/@robots-that-learn}}\\
官方 playlist：\href{https://www.youtube.com/playlist?list=PLPaC96j0xdLcYLTSoSk9PO1Yg-1udJd-S}{\nolinkurl{https://www.youtube.com/playlist?list=PLPaC96j0xdLcYLTSoSk9PO1Yg-1udJd-S}}

该官方频道当前公开可见的课程 playlist 只有这一条，公开 playlist 共包含 15 个视频。
这与课程页和课程 notes 的讲次切分并不完全一致，因此本讲义在整书层明确保留以下 coverage 说明：
\begin{itemize}[leftmargin=2em]
\item Lecture 1A: Introduction：当前 run 未发现对应的公开 YouTube 视频，章节内容主要由官方 course notes、slides 与课程页材料重建。
\item Lecture 1B: Biomechanics of Walking and Running：当前 run 未发现对应的公开 YouTube 视频，章节内容主要由官方 course notes、slides 与课程页材料重建。
\item Lecture 2B: The Human Hand and Dexterous Object Manipulation：当前 run 未发现对应的公开 YouTube 视频，章节内容主要由官方 course notes、slides 与课程页材料重建。
\end{itemize}

\section*{课程目录}
\addcontentsline{toc}{section}{课程目录}
\begin{longtable}{p{0.08\textwidth}p{0.60\textwidth}p{0.14\textwidth}p{0.10\textwidth}}
\textbf{讲次} & \textbf{主题} & \textbf{日期} & \textbf{Schedule}\\
\hline
01 & Lecture 1A: Introduction & 2024-09-09 & 1A\\
02 & Lecture 1B: Biomechanics of Walking and Running & 2024-09-09 & 1B\\
03 & Lecture 2A: Robot Mechanisms - Kinematics and Dynamics & 2024-09-16 & 2A\\
04 & Lecture 2B: The Human Hand and Dexterous Object Manipulation & 2024-09-16 & 2B\\
05 & Lecture 3: Robot Hands; Proprioception and Tactile Perception & 2024-09-23 & 3AB\\
06 & Lecture 4A: Vision for Action & 2024-09-30 & 4A\\
07 & Lecture 4B: The Developmental Perspective on Motor Control & 2024-09-30 & 4B\\
08 & Lecture 5A: Robot Dynamics, Control, and Motion Planning & 2024-10-07 & 5A\\
09 & Lecture 5B: Computational Neuroscience Perspective on Prediction and Control & 2024-10-07 & 5B\\
10 & Lecture 6A: Reinforcement Learning & 2024-10-14 & 6A\\
11 & Lecture 6B: Reinforcement Learning & 2024-10-14 & 6B\\
12 & Lecture 7: Behavior Cloning & 2024-10-21 & 7AB\\
13 & Lecture 8: Visual Imitation & 2024-10-28 & 8AB\\
14 & Lecture 9: Case Studies in Locomotion & 2024-11-04 & 9AB\\
15 & Lecture 10: Case Studies in Navigation & 2024-11-18 & 10AB\\
16 & Lecture 11A: Case Studies in Dexterous Manipulation & 2024-11-25 & 11A\\
17 & Lecture 11B: Case Studies in Dexterous Manipulation & 2024-11-25 & 11B\\
18 & Lecture 12: Long Horizon Planning and the Role of Language & 2024-12-02 & 12AB\\
\end{longtable}

\section{Lecture 1A: Introduction}
\includepdf[pages=-,pagecommand={\thispagestyle{plain}}]{../lectures/01_introduction/lecture_01_note.pdf}

\section{Lecture 1B: Biomechanics of Walking and Running}
\includepdf[pages=-,pagecommand={\thispagestyle{plain}}]{../lectures/02_biomechanics_of_walking_and_running/lecture_02_note.pdf}

\section{Lecture 2A: Robot Mechanisms - Kinematics and Dynamics}
\includepdf[pages=-,pagecommand={\thispagestyle{plain}}]{../lectures/03_robot_kinematics_and_dynamics/lecture_03_note.pdf}

\section{Lecture 2B: The Human Hand and Dexterous Object Manipulation}
\includepdf[pages=-,pagecommand={\thispagestyle{plain}}]{../lectures/04_human_hand_and_dexterous_object_manipulation/lecture_04_note.pdf}

\section{Lecture 3: Robot Hands; Proprioception and Tactile Perception}
\includepdf[pages=-,pagecommand={\thispagestyle{plain}}]{../lectures/05_robot_hands_proprioception_and_touch/lecture_05_note.pdf}

\section{Lecture 4A: Vision for Action}
\includepdf[pages=-,pagecommand={\thispagestyle{plain}}]{../lectures/06_vision_for_action/lecture_06_note.pdf}

\section{Lecture 4B: The Developmental Perspective on Motor Control}
\includepdf[pages=-,pagecommand={\thispagestyle{plain}}]{../lectures/07_developmental_perspective_on_motor_control/lecture_07_note.pdf}

\section{Lecture 5A: Robot Dynamics, Control, and Motion Planning}
\includepdf[pages=-,pagecommand={\thispagestyle{plain}}]{../lectures/08_robot_dynamics_control_and_motion_planning/lecture_08_note.pdf}

\section{Lecture 5B: Computational Neuroscience Perspective on Prediction and Control}
\includepdf[pages=-,pagecommand={\thispagestyle{plain}}]{../lectures/09_prediction_and_control_from_computational_neuroscience/lecture_09_note.pdf}

\section{Lecture 6A: Reinforcement Learning}
\includepdf[pages=-,pagecommand={\thispagestyle{plain}}]{../lectures/10_reinforcement_learning_part_a/lecture_10_note.pdf}

\section{Lecture 6B: Reinforcement Learning}
\includepdf[pages=-,pagecommand={\thispagestyle{plain}}]{../lectures/11_reinforcement_learning_part_b/lecture_11_note.pdf}

\section{Lecture 7: Behavior Cloning}
\includepdf[pages=-,pagecommand={\thispagestyle{plain}}]{../lectures/12_behavior_cloning/lecture_12_note.pdf}

\section{Lecture 8: Visual Imitation}
\includepdf[pages=-,pagecommand={\thispagestyle{plain}}]{../lectures/13_visual_imitation/lecture_13_note.pdf}

\section{Lecture 9: Case Studies in Locomotion}
\includepdf[pages=-,pagecommand={\thispagestyle{plain}}]{../lectures/14_case_studies_in_locomotion/lecture_14_note.pdf}

\section{Lecture 10: Case Studies in Navigation}
\includepdf[pages=-,pagecommand={\thispagestyle{plain}}]{../lectures/15_case_studies_in_navigation/lecture_15_note.pdf}

\section{Lecture 11A: Case Studies in Dexterous Manipulation}
\includepdf[pages=-,pagecommand={\thispagestyle{plain}}]{../lectures/16_dexterous_manipulation_part_a/lecture_16_note.pdf}

\section{Lecture 11B: Case Studies in Dexterous Manipulation}
\includepdf[pages=-,pagecommand={\thispagestyle{plain}}]{../lectures/17_dexterous_manipulation_part_b/lecture_17_note.pdf}

\section{Lecture 12: Long Horizon Planning and the Role of Language}
\includepdf[pages=-,pagecommand={\thispagestyle{plain}}]{../lectures/18_long_horizon_planning_and_language/lecture_18_note.pdf}

\appendix
\section*{补充附录}
\addcontentsline{toc}{section}{补充附录}
以下章节用于把 Spring 2026 的最新课程更新，以及 MIT Underactuated Spring 2024 的理论骨架，纳入当前主书的系统化补充。

\section{Appendix A: Spring 2026 最新课程主线与研究重排}
\includepdf[pages=-,pagecommand={\thispagestyle{plain}}]{../lectures/19_spring2026_latest_course_updates/lecture_19_note.pdf}

\section{Appendix B: Spring 2026 的 Diffusion Models 与 Video World Models 补充}
\includepdf[pages=-,pagecommand={\thispagestyle{plain}}]{../lectures/20_spring2026_diffusion_and_world_models/lecture_20_note.pdf}

\section{Appendix C: Underactuated 课程的动力学、控制与规划基础}
\includepdf[pages=-,pagecommand={\thispagestyle{plain}}]{../lectures/21_underactuated_dynamics_control_and_planning_foundations/lecture_21_note.pdf}

\section{Appendix D: Underactuated 课程的 locomotion、contact 与 robustness 补充}

\includepdf[pages=-,pagecommand={\thispagestyle{plain}}]{../lectures/22_underactuated_locomotion_contact_and_robustness/lecture_22_note.pdf}

\end{document}
